% !TeX program = xelatex
% 从本文件所在目录编译。完整主题与 assets 随课件提供。
% 主讲1--21页；集合题22--28页；第29页结束。研究信息核查至2026-09-11。
\documentclass[aspectratio=169,10pt]{beamer}
\usepackage[UTF8,fontset=fandol]{ctex}
\IfFontExistsTF{Microsoft YaHei}{\setCJKsansfont{Microsoft YaHei}}{}
\usetheme{GaoZhongYuan}
\usepackage{booktabs,array}
\setbeamerfont{normal text}{size=\fontsize{12}{17}\selectfont}
\setbeamerfont{itemize/enumerate body}{size=\fontsize{12}{17}\selectfont}
\setbeamerfont{itemize/enumerate subbody}{size=\fontsize{11}{15}\selectfont}
\setbeamertemplate{itemize items}[circle]
\newcommand{\key}[1]{{\color{campusgreen}\bfseries #1}}
% 出处保留在源文件与授课说明，放映版不显示。
\newcommand{\source}[2]{\par\vfill{\fontsize{7}{9}\selectfont\strut\par}}
\newcommand{\smallbody}{\fontsize{11}{15}\selectfont}
\newcommand{\backtosummary}{\hfill\hyperlink{summary}{\beamerbutton{回到课堂小结}}}
\hypersetup{pdftitle={AI伴学 第一课},pdfauthor={华光辉},unicode=true}


\title[AI伴学 第一课]{\texorpdfstring{AI伴学\quad 第一课}{AI伴学 第一课}}
\author{华光辉}
\date{2026.09.12}
\begin{document}
\usebeamerfont{normal text}
\campusmaketitle

\begin{frame}[t]{这门课怎么上}
  这学期，我们试着用AI帮助自己读书、思考和解决数学问题。
  \par\vspace{8mm}
  \begin{tabular}{@{}lp{109mm}@{}}
    \key{课程要求} & 期末交什么，怎样认定学分。\\[5mm]
    \key{认识AI} & 看两段视频，聊一聊它的发展和数学研究中的新进展。\\[5mm]
    \key{现场演示} & 我怎样使用自己搭建的Scholar系统。\\
  \end{tabular}
  \vfill
  如果时间够，我们再看一道集合题。
\end{frame}


\begin{frame}[t,label=outlineone]{这学期学什么（一）}
  \smallbody
  \key{前10课：用几个新的数学主题，练习怎样借助AI学习。}
  \par\vspace{3mm}
  \begin{tabular}{@{}lp{113mm}@{}}
    \key{第1—2课} & \textbf{认识AI，配置自己的学习助手}\\[1mm]
                   & 了解AI的发展与数学应用；确定目标，配置并测试助手。\\[2mm]
    \key{第3—5课} & \textbf{复数与平面}\\[1mm]
                   & 从$x^2+1=0$出发，学习复数运算，探索旋转与伸缩。\\[2mm]
    \key{第6—7课} & \textbf{图论与一笔画}\\[1mm]
                   & 认识点、边和度数；判断一笔画，把路线问题画成图。\\[2mm]
    \key{第8—10课} & \textbf{迭代与分形}\\[1mm]
                   & 反复代入、重复作图，观察稳定、循环和自相似，交流发现。\\
  \end{tabular}
  \vfill
  {\fontsize{9}{12}\selectfont 全学期暂定18课时，每课时40分钟。}
\end{frame}

\begin{frame}[t,label=outlinetwo]{这学期学什么（二）}
  \smallbody
  \key{后8课：把方法用回课内，看看自己是否真正学会。}
  \par\vspace{3mm}
  \begin{tabular}{@{}lp{113mm}@{}}
    \key{第11课} & \textbf{迁移到集合、函数等课内内容}\\[1mm]
                  & 用例子、反例和自己的解释，检查概念是否理解。\\[2mm]
    \key{第12—15课} & \textbf{练习、纠错，再独立检验}\\[1mm]
                     & 独立作答、分析错因、针对性练习，再隔几天做新同类题。\\[2mm]
    \key{第16课} & \textbf{高考能力衔接}\\[1mm]
                  & 用已经学过的知识，练习条件转化、分类讨论与规范表达。\\[2mm]
    \key{第17—18课} & \textbf{成果展示与独立检验}\\[1mm]
                     & 展示学习过程，完成分层任务，为自选新主题设计学习计划。\\
  \end{tabular}
  \vfill
  {\fontsize{9}{12}\selectfont 教学内容暂定；课内题目随实际教学进度调整。}
\end{frame}

\begin{frame}[t]{期末评价与学分}
  \renewcommand{\arraystretch}{1.5}
  \begin{tabular}{@{}llp{72mm}@{}}
    \key{40\%} & 方法与反思档案 & 保存自己的尝试、AI的建议和修改过程。\\
    \key{30\%} & 独立数学任务 & 不借助AI，解释或完成一道数学任务。\\
    \key{30\%} & 探究与迁移展示 & 讲清一个问题，以及这次学到的方法。\\
  \end{tabular}\par
  \vspace{7mm}
  平时把学习记录存好，期末从中选出能说明自己进步的材料。
  \vfill
  {\smallbody 学分数及认定条件以学校公布为准，目前课程申报表中尚待核定。}
\end{frame}

\begin{frame}[t]{学习记录可以这样写}
  \begin{tabular}{@{}lp{114mm}@{}}
    \key{题目} & 我这次想解决什么？\\[5mm]
    \key{尝试} & 我已经做到了哪一步？\\[5mm]
    \key{求助} & 我问了什么？哪条建议有用，哪条不对？\\[5mm]
    \key{回头再做} & 不看对话记录，我能不能自己完成？\\
  \end{tabular}
  \vfill
  一份记录不用很长，但要看得出自己的思考。
\end{frame}

\begin{frame}[t]{人工智能并不是这几年才出现的}
  \key{1956年\quad 达特茅斯夏季研究计划}
  \par\vspace{3mm}
  研究者提出：机器能否使用语言、形成概念、解决问题？
  \par\vspace{7mm}
  早期的一条思路，是把知识和推理规则写给计算机。
  \[
    \text{已知条件}\quad +\quad \text{推理规则}\quad \longrightarrow\quad \text{结论}
  \]
  \vfill
  可是，认出一张照片里的猫，能写成几条简单规则吗？
\end{frame}

\begin{frame}[t]{让计算机从例子中学习}
  给出许多例子，调整模型，使它逐渐学会识别和预测。
  \par\vspace{7mm}
  \begin{tabular}{@{}lp{117mm}@{}}
    \key{2012} & AlexNet在图像识别比赛中取得突出成绩。\\[6mm]
    \key{2016} & AlphaGo以4比1战胜李世石。\\[2mm]
              & 它结合了神经网络、搜索和强化学习。\\
  \end{tabular}
  \vfill
  棋盘上的变化太多，棋手和计算机都需要判断哪些变化值得继续看。
\end{frame}

\begin{frame}[t,label=neuron]{一个“神经元”怎样计算}
  \smallbody
  把图片等信息变成数。每个输入乘一个系数，再加起来。
  \par\vspace{4mm}
  \begin{center}
  \begin{tikzpicture}[x=1cm,y=1cm,>=stealth]
    \node (a) at (0,.65) {$x_1=1$};
    \node (b) at (0,-.65) {$x_2=0.5$};
    \node[draw,circle,minimum size=11mm] (sum) at (3,0) {$s$};
    \node[draw,minimum width=29mm,minimum height=10mm] (act) at (6,0) {负数变成0};
    \node (out) at (9,0) {输出$0.3$};
    \draw[->] (a)--node[above,sloped] {$\times\,0.6$}(sum);
    \draw[->] (b)--node[below,sloped] {$\times\,0.4$}(sum);
    \node[below] at (3,-.65) {再减去$0.5$};
    \draw[->] (sum)--(act);\draw[->] (act)--(out);
  \end{tikzpicture}
  \end{center}
  \[
    s=\underbrace{0.6}_{\text{权重}}x_1+\underbrace{0.4}_{\text{权重}}x_2
      \underbrace{-0.5}_{\text{偏置}}=0.3,
    \qquad y=\max\{0,s\}.
  \]
  \vfill
  最后这一步叫\key{激活函数}。这里用的是ReLU：负数变成0，非负数保留。
  \par\vspace{2mm}
  {\fontsize{9}{12}\selectfont 数字仅作课堂示例；这个输出还不是识别结论，也不是概率。}
\end{frame}

\begin{frame}[t,label=network]{连成网络，再用例子训练}
  \begin{center}
  \begin{tikzpicture}[x=1cm,y=1cm,>=stealth]
    \foreach \i/\y in {1/.6,2/-.6}{\node[draw,circle,minimum size=6mm] (i\i) at (0,\y) {};}
    \foreach \j/\y in {1/.9,2/0,3/-.9}{\node[draw,circle,minimum size=6mm] (h\j) at (3.6,\y) {};}
    \node[draw,circle,minimum size=7mm] (o) at (7.2,0) {};
    \foreach \i in {1,2}{\foreach \j in {1,2,3}{\draw[->,gray] (i\i)--(h\j);}}
    \foreach \j in {1,2,3}{\draw[->,gray] (h\j)--(o);}
    \node[above] at (0,1.2) {输入};
    \node[above] at (3.6,1.2) {中间一层};
    \node[above] at (7.2,1.2) {输出};
  \end{tikzpicture}
  \end{center}
  \smallbody
  中间层的圆点做这样的计算，再把结果传给下一层。
  \par\vspace{4mm}
  \key{训练时}\quad 输入例子\ $\longrightarrow$\ 算出结果\ $\longrightarrow$\ 和已知答案比较。
  \par\vspace{3mm}
  根据误差调整权重和偏置，再试一次，让\textbf{总体误差}尽量减小。
  \vfill
  训练完，还要用\textbf{没参与训练的例子}检查效果。
\end{frame}

\begin{frame}[t,label=alphagovideo]{纪录片《AlphaGo》预告}
  \vspace{4mm}
  先看一段预告，感受一下当时人们怎样看待这场比赛。
  \par\vspace{12mm}
  \begin{center}
    \href{https://www.bilibili.com/video/BV13x411x7kY/}{\color{campusgreen}\bfseries\fontsize{19}{25}\selectfont\underline{播放视频}}
    \par\vspace{4mm}
    中文字幕\qquad 1分30秒
  \end{center}
  \vfill
  {\smallbody 点击后在浏览器播放，需要联网。}
\end{frame}

\begin{frame}[t,label=hintonvideo]{Hinton谈自己的研究经历}
  \vspace{4mm}
  一项今天看来很成功的研究，在当时未必受到认可。
  \par\vspace{12mm}
  \begin{center}
    \href{https://www.bilibili.com/video/BV1xr4y1j7hR/}{\color{campusgreen}\bfseries\fontsize{19}{25}\selectfont\underline{播放访谈}}
    \par\vspace{4mm}
    10分27秒
  \end{center}
  \vfill
  {\smallbody 点击后在浏览器播放，需要联网。}
\end{frame}

\begin{frame}[t]{今天的AI为什么能和我们对话}
  \key{2017年\quad Transformer}
  \par\vspace{3mm}
  注意力机制帮助模型处理上下文中不同信息的联系。
  \par\vspace{7mm}
  大语言模型能接着上下文生成文字，也能写代码、解释问题。
  \par\vspace{7mm}
  接上检索、计算和证明检查工具之后，能做的事情更多了。
  \vfill
  比如，一道题的答案看起来很像标准答案，它就一定对吗？
\end{frame}

% 文献与核查说明见《前沿数学核查记录.md》，放映版不显示出处。

\begin{frame}[t]{从解竞赛题，到研究新问题}
  \key{2025年\quad 国际数学奥林匹克（IMO）试题}
  \par\vspace{8mm}
  增强版Gemini Deep Think解出6道题中的5道，\par
  得到35分（满分42分），达到当年的金牌标准。
  \par\vspace{10mm}
  竞赛题有已知的答案和评分标准。\par
  研究问题，则可能连数学家也还不知道答案。
  \vfill
  下面从大家熟悉的水流，说起一个研究问题。
\end{frame}

\begin{frame}[t,label=nsintro]{NS方程：用数学描述水怎样流动}
  搅动一杯水，水会旋转；停止搅动后，它又会慢慢平静下来。
  \par\vspace{6mm}
  \key{基本想法：把牛顿第二定律 $F=ma$ 用到流体上。}
  \par\vspace{5mm}
  \begin{tabular}{@{}lp{113mm}@{}}
    压力差 & 水的不同位置，受到的推挤不同。\\[4mm]
    黏性 & 相邻的流体彼此牵制，倾向于减小速度差。\\[4mm]
    外力 & 例如重力，或搅动带来的作用。\\
  \end{tabular}
  \vfill
  NS（Navier--Stokes）方程把这些作用联系起来，\par
  描述各处的流速、方向怎样随时间变化。
\end{frame}

\begin{frame}[t,label=nsquestion]{难在哪里：一直算下去，会发生什么}
  \smallbody
  从光滑的初始状态出发，流动能一直保持光滑吗？\par
  还是会在\key{有限时间}内，出现越来越尖锐的变化？
  \par\vspace{4mm}
  先看一个熟悉的函数：$y=\dfrac{1}{1-t}$，其中$0\le t<1$。
  \par\vspace{3mm}
  \begin{center}
  \renewcommand{\arraystretch}{1.35}
  \setlength{\tabcolsep}{7mm}
  \begin{tabular}{c|cccc}
    $t$ & $0$ & $0.9$ & $0.99$ & $0.999$\\\hline
    $y$ & $1$ & $10$ & $100$ & $1000$
  \end{tabular}
  \end{center}
  \par\vspace{3mm}
  $t$越来越接近1，$y$却没有上界。\par
  “有限时间爆破”说的是这种数学上的无界增长。
  \vfill
  {\fontsize{9}{12}\selectfont 这个函数只是类比，不是NS方程的解；“爆破”也不是水真的爆炸。}
\end{frame}

\begin{frame}[t,label=nsprogress]{AI参与了什么}
  \smallbody
  \key{2026年9月8日\quad 一项新公布的研究}
  \par\vspace{4mm}
  OpenAI报告：AI系统构造了一种特殊的三维流动，\par
  在\textbf{光滑外力}作用下，数学模型中的速度可在有限时间内无界增长。
  \par\vspace{5mm}
  研究团队公开了证明论文，也公开了供计算机核查推理的代码。
  \par\vspace{5mm}
  对数学研究而言，AI正在参与寻找新构造和证明。
  \vfill
  {\fontsize{10}{14}\selectfont 结果刚公布，仍需进一步独立审查。\par
  这也不意味着现实中的水会达到无穷大的速度。}
\end{frame}

\begin{frame}[t]{我怎样使用Scholar学习}
  我把教材、笔记和学习记录放在同一个系统里。
  \par\vspace{7mm}
  \begin{tabular}{@{}lp{109mm}@{}}
    \key{读资料} & 打开教材，找到自己没看懂的一句话。\\[5mm]
    \key{问问题} & 请AI解释，或者检查自己的理解。\\[5mm]
    \key{做练习} & 先写尝试，需要时再要一个提示。\\[5mm]
    \key{过几天再看} & 回顾记录，检查自己还能不能完成。\\
  \end{tabular}
  \vfill
  下面切换到系统，实际做一遍。
\end{frame}

\begin{frame}[t]{演示：读到一句没看懂的话}
  在Library中打开一份资料，选中一个定义或证明步骤。
  \par\vspace{8mm}
  \key{我会这样问}
  \par\vspace{3mm}
  “这一步为什么成立？用了前面的哪一个条件？\par
  请先解释这一小步，再给一个简单例子。”
  \par\vspace{7mm}
  看完回答，再回到原文，对照它解释的那一步。
  \vfill
  Ask \& Understand：适合读书时的解释和追问。
\end{frame}

\begin{frame}[t]{演示：一道题暂时做不下去}
  \begin{enumerate}\setlength{\itemsep}{5mm}
    \item 先把已经写出的思路保存下来。
    \item 点Request one hint，只要一个提示。
    \item 自己继续写，说明改动了哪一步。
    \item 安排1天后、7天后的复习。
  \end{enumerate}
  \vfill
  Practice \& Prove：适合练习和证明。
  \par\vspace{3mm}
  \hyperlink{setproblem}{\beamerbutton{集合题}}\hfill
  \hyperlink{summary}{\beamerbutton{课堂小结}}
\end{frame}

\begin{frame}[t,label=summary]{下课前}
  \vspace{4mm}
  写下一个你希望这学期弄懂的问题。
  \par\vspace{9mm}
  可以是一道题、一本书里的一句话，\par
  也可以是今天听到的一个数学问题。
  \par\vspace{9mm}
  下次课，我们从自己的问题开始，配置和试用学习助手。
  \vfill
  \hyperlink{setproblem}{\beamerbutton{集合题选讲}}\hfill
  \hyperlink{finish}{\beamerbutton{结束}}
\end{frame}
\appendix
\begin{frame}[t,label=setproblem]{选讲：集合中的一个难题}
  \smallbody
  已知集合 $A=\{a_1,a_2,\ldots,a_n\}$，其中 $n\in\mathbb N^*$ 且 $n\ge4$，
  $a_i\in\mathbb N^*\ (i=1,2,\ldots,n)$。非空集合 $B\subseteq A$，
  记 $T(B)$ 为集合 $B$ 中所有元素之和；若 $B=\{b\}$，则 $T(B)=b$。
  \par\vspace{3mm}
  （1）若 $A=\{1,2,5,6,7,8\}$，$T(B)=8$，写出所有可能的集合 $B$。
  \par\vspace{3mm}
  （2）若 $A=\{3,4,5,9,10,11\}$，$B=\{b_1,b_2,b_3\}$，且 $T(B)$ 是
  $12$ 的倍数，求集合 $B$ 的个数。
  \par\vspace{3mm}
  （3）若 $a_i\in\{1,2,3,\ldots,2n-1\}\ (i=1,2,\ldots,n)$，证明：
  存在非空集合 $B\subseteq A$，使得 $T(B)$ 是 $2n$ 的倍数。
  \par\vfill
  \backtosummary
\end{frame}

\begin{frame}[t]{选讲：先让AI帮哪一步}
  \smallbody
  \key{先读题}
  \par “请复述题目。$T(B)$ 是元素之和，$B$ 非空，集合元素互异。\par
  先检查条件，不要直接给完整证明。”
  \par\vspace{3mm}
  \key{再找思路}
  \par “前两问怎样分类才能不重不漏？第三问能否先看两个数的和？”
  \par\vspace{3mm}
  \key{最后检查证明}
  \par “如果按和为 $2n$ 配对，单独的 $n$ 怎么处理？\par
  每一步推理使用了什么条件？请给我一个小提示。”
  \vfill
   列举是否完整？找到了例子，还是证明了一般结论？
  \par\vspace{2mm}\backtosummary
\end{frame}

\begin{frame}[t]{选讲：第一问的完整列举}
  \key{$A=\{1,2,5,6,7,8\}$，元素和为8}
  \par\vspace{4mm}
  \begin{tabular}{@{}lp{103mm}@{}}
    1个元素 & $\{8\}$\\[3mm]
    2个元素 & $\{1,7\},\ \{2,6\}$\\[3mm]
    3个元素 & $\{1,2,5\}$\\[3mm]
    至少4个元素 & 最小的和为 $1+2+5+6=14>8$，不可能。\\
  \end{tabular}
  \vfill
  \key{共4个。}\quad 按元素个数分类，可以说明为什么没有遗漏。
  \par\vspace{3mm}\backtosummary
\end{frame}

\begin{frame}[t]{选讲：第二问的完整列举}
  \key{$A=\{3,4,5,9,10,11\}$，取3个元素}
  \par\vspace{3mm}
  最小和为 $3+4+5=12$，最大和为 $9+10+11=30$。\par
  因此只需检查 $T(B)=12$ 或 $T(B)=24$。
  \par\vspace{4mm}
  \begin{tabular}{@{}lp{109mm}@{}}
    $T(B)=12$ & $\{3,4,5\}$\\[4mm]
    $T(B)=24$ & $\{3,10,11\},\ \{4,9,11\},\ \{5,9,10\}$\\
  \end{tabular}
  \vfill
  \key{共4个。}\quad 集合不计顺序，同一组元素只算一次。
  \par\vspace{3mm}\backtosummary
\end{frame}

\begin{frame}[t]{选讲：配对法的关键分支}
  \smallbody
  将 $1,2,\ldots,2n-1$ 分为
  \[
    \{1,2n-1\},\ \{2,2n-2\},\ \ldots,\ \{n-1,n+1\},\ \{n\}.
  \]
  \key{情况一}\quad $A$ 取到了某一对的两个数。
  \par 这两个数的和就是 $2n$，取它们组成 $B$ 即可。
  \par\vspace{4mm}
  \key{情况二}\quad 每一对至多取一个数。
  \par 因为 $|A|=n$，所以必须有 $n\in A$，且每一对恰好取一个数。
  \par\vspace{4mm}
  \key{注意}
  \par 这里共有 $n$ 组，选 $n$ 个数，不保证某一组一定选到两个。
  \vfill\backtosummary
\end{frame}

\begin{frame}[t]{选讲：一个余数引理}
  \smallbody
  \key{引理}\quad 若 $n-1$ 个整数的任意非空选取之和都不被 $n$ 整除，\par
  则这些整数除以 $n$ 的余数全部相同。
  \par\vspace{3mm}
  \key{证明}\quad 将它们任意排列为 $c_1,c_2,\ldots,c_{n-1}$，令
  \[
     S_1=c_1,\quad S_2=c_1+c_2,\quad\ldots,\quad S_{n-1}=c_1+\cdots+c_{n-1}.
  \]
  各 $S_i$ 的余数非零；若 $S_i,S_j\ (i<j)$ 同余，\par
  则 $c_{i+1}+\cdots+c_j$ 被 $n$ 整除，矛盾。
  \par\vspace{3mm}
  所以这 $n-1$ 个余数恰好是 $1,2,\ldots,n-1$。
  \par 交换 $c_1,c_2$ 后，只有第一个部分和改变，其余不变，\par
  故 $c_1\equiv c_2\pmod n$。任意两项都可放在前两位，结论成立。
  \vfill\backtosummary
\end{frame}

\begin{frame}[t]{选讲：完成第三问的证明}
  \smallbody
  在情况二中，反设没有非空子集的元素和被 $2n$ 整除。\par
  令 $C=A\setminus\{n\}$，则 $|C|=n-1$。
  \par\vspace{3mm}
  \key{第一步}\quad $C$ 的任意非空子集的和都不被 $n$ 整除。
  \par 否则设其和为 $kn$：若 $k$ 为偶数，已经得到结论；\par
  若 $k$ 为奇数，再加入元素 $n$，和为 $(k+1)n$，也被 $2n$ 整除。
  \par\vspace{3mm}
  \key{第二步}\quad 由引理，$C$ 的所有元素模 $n$ 同余。
  \par 在 $1,2,\ldots,2n-1$ 中，同一非零余数最多对应两个数：$r,r+n$。
  \par 但 $|C|=n-1\ge3$，且元素互异，矛盾。
  \par\vspace{3mm}
  因此一定存在所求的非空集合 $B$。\hfill$\square$
  \vfill\backtosummary
\end{frame}


\hypertarget{finish}{}
\campusclosing[下次课见]
\end{document}
